\subsection{Soft Actor-Critic (SAC) Foundation}

We select Soft Actor-Critic over alternative RL algorithms for three fundamental reasons. First, sample efficiency via off-policy learning is critical because real-world ASV deployments are expensive and limited. SAC maintains a replay buffer enabling efficient reuse of past transitions, achieving strong performance with 5--10 times fewer environment interactions compared to on-policy methods like PPO. Second, continuous action spaces are essential for smooth motor control. ASV thrust, steering angle, and sensor duty cycles are inherently continuous; discretization approaches (DQN with action binning) introduce quantization error and explosion in action space dimensionality. Third, entropy regularization prevents premature convergence in maritime uncertainty where sensor noise, current variability, and incomplete observability reward stochastic exploration over exploitation.

In contrast, deterministic alternatives such as TD3 sacrifice exploration benefits critical for discovering energy-efficient strategies in unseen water conditions. PPO and other on-policy methods consume vastly more real-world data. DQN variants suffer from discretization penalties and combinatorial action explosion.

SAC maintains three network families. A stochastic policy $\pi(a|s)$ maps state to action distribution. Dual critics $Q_1(s,a)$ and $Q_2(s,a)$ independently estimate action-value functions. Target networks $\tilde{Q}_{\text{target}}(s,a)$ stabilize learning. The optimization objective combines policy gradient with entropy regularization:
\begin{align}
J(\pi) = \mathbb{E}_{s \sim \mathcal{D}}[\mathbb{E}_{a \sim \pi}[-\alpha \log \pi(a|s) + Q_{\min}(s,a)]]
\end{align}
where $\alpha$ is a learned temperature coefficient and $Q_{\min} = \min(Q_1, Q_2)$ prevents overestimation bias that could lead to unsafe energy-depleting actions.

\subsection{State Space Design (87D Vector)}

The state vector $s_t$ encodes multimodal information enabling mission-critical decisions while remaining computationally lightweight for real-time 10 Hz control loops.

Position and dynamics (9D) include Cartesian coordinates $(x, y)$, heading $\theta$, linear velocities $(v_x, v_y)$, angular velocity $\omega$, linear acceleration, angular acceleration, and distance to home base. These elements ground the policy in absolute space and provide trajectory context.

Battery state (6D) encompasses state-of-charge (SoC) as a percentage in $[0, 1]$, absolute energy remaining in kWh, estimated energy required to return home given current distance and environmental resistance, energy buffer safety margin indicating how much above the safe return threshold, battery degradation factor (multiplicative decay in capacity), and cycling count (number of charge cycles affecting long-term capacity). This representation enables the policy to distinguish between a critically-low 5\% SoC versus a more recoverable 10\% SoC with different remaining distances.

Environmental conditions (12D) provide the policy with hydrodynamic and atmospheric context. Water current is encoded as 2D velocity (east/north components). Wave state includes significant wave height and period. Wind is encoded as 2D velocity vector. Water depth (affecting drag and sensor performance), temperature (affecting battery and sensor behavior), and salinity (affecting sensor calibration) provide physical properties. GPS quality indicator (0--1 scale) captures localization uncertainty. Time-of-day normalized to $[0, 1]$ enables diurnal adaptation. Weather is one-hot encoded as clear, rain, or fog.

Water-body type classification (4D) is encoded as one-hot vector from the set $\{\text{lake}, \text{river}, \text{coastal\_ocean}, \text{open\_ocean}\}$. This is inferred from environmental signatures during training (e.g., steady current patterns suggest river, wave periodicity suggests ocean). During deployment, this can be initialized from mission metadata or learned via an auxiliary classifier.

Mission progress (18D) includes coverage percentage as $[0, 1]$. Frontier density is represented as an 8×8 occupancy grid compressed to 10D via learned embedding (or used raw if computational budget permits). Distance to nearest frontier in meters. Number of active (unexplored) frontier clusters. Average frontier information value (proxy for unexplored region richness). Cumulative information gain (integrated sensor measurements). Mission elapsed time normalized to mission budget. Estimated depletion time at current consumption rate.

Sensor states (16D) capture per-sensor activation status (4D binary vector for Radar, Lidar, Camera, Water quality). Duty cycle history (4D) records the fractional activation time for each sensor over the last timestep. Information gain per sensor (4D) estimates each sensor's contribution to coverage/knowledge in the current location. Last activation timestamp per sensor (4D) enables the policy to enforce minimum inter-activation intervals.

Historical context (22D) remembers the last five consumption rates (5D), last five coverage gains (5D), last five actions in terms of thrust and steering (5 actions × 2 = 10D), frontier revisit count (current and total, 2D). History enables the policy to detect trends and adapt smoothly.

This 87-dimensional representation balances information richness with computational feasibility. State construction on commodity CPU hardware requires approximately 2--3ms, enabling 10 Hz control loops with 10ms timesteps.

\subsection{Hierarchical Action Space}

To balance precision with interpretability and computational efficiency, we employ a two-stage hierarchical action space.

Stage 1 (Discrete, 4 choices) selects the motion mode. Exploration mode ($a_1 = 0$) activates aggressive coverage with maximum feasible thrust and full sensor activation. Energy-saving mode ($a_1 = 1$) minimizes thrust, activates only one or two sensors, and may exploit current-assisted drift to reduce energy consumption. Return-to-base mode ($a_1 = 2$) overrides all other objectives and plots a direct path toward home base with minimal sensing, ensuring timely return before critical battery depletion. Information-gathering mode ($a_1 = 3$) maintains stationary or near-stationary position while activating all sensors to maximize local measurement density.

Stage 2 (Continuous, 7D per mode) parameterizes control within the selected mode. Motion commands include thrust percentage in $[-1.0, +1.0]$ (negative for reverse), steering angle in $[-\pi/2, +\pi/2]$ radians, and speed target in $[0, v_{\max}]$ m/s. Sensor duty cycles encode probabilistic activation for each of four sensors in $[0, 1]$, where the probability is sampled stochastically at each control step, enabling intermittent sensing that averages to a learned duty cycle over longer timescales.

This hierarchical decomposition reduces effective action space from approximately $10^7$ continuous combinations to $4 \times 7 = 28$ effective dimensions while maintaining expressiveness and interpretability. The discrete mode selection is easily logged for human operators, providing transparency into mission strategy.

\subsection{Composite Reward Function}

The reward at timestep $t$ combines multiple objectives:
\begin{align}
r_t = \lambda_{\text{cov}} r^{\text{cov}}_t + \lambda_{\text{info}} r^{\text{info}}_t + \lambda_{\text{energy}} r^{\text{energy}}_t + \lambda_{\text{safety}} r^{\text{safety}}_t + \lambda_{\text{eff}} r^{\text{eff}}_t
\end{align}

Coverage reward ($\lambda_{\text{cov}} = 1.0$) is:
\begin{align}
r^{\text{cov}}_t = \frac{\text{new\_area\_explored}}{\text{max\_area\_per\_step}}
\end{align}
normalized to $[0, 1]$ by the theoretical maximum coverage achievable per timestep given sensor footprints and vehicle velocity.

Information reward ($\lambda_{\text{info}} = 0.7$) aggregates per-sensor contributions:
\begin{align}
r^{\text{info}}_t = \sum_{i=1}^{4} \mathbb{1}[\text{sensor}_i \text{ active}] \times \text{info\_gain}_i \times \text{quality\_factor}_i
\end{align}
where information gain depends on sensor type. Radar measures object detections in unexplored regions. Lidar provides 3D structure for bathymetric mapping. Camera captures visual features (edge density, texture). Water quality sensor measures chemical readings in under-sampled locations. Quality factors degrade with environmental conditions (e.g., camera quality factor approaches 0 in dense fog; Lidar quality factor decreases in rough seas).

Energy efficiency penalty ($\lambda_{\text{energy}} = -0.8$) discourages wasteful consumption:
\begin{align}
r^{\text{energy}}_t = -\frac{\text{energy\_consumed}\_t}{\text{baseline\_energy\_per\_step}}
\end{align}
normalized such that station-keeping (minimum energy) gives penalty near zero.

Safety reward ($\lambda_{\text{safety}} = -2.0$, severe penalty) enforces energy constraints:
\begin{align}
r^{\text{safety}}_t = \begin{cases}
0 & \text{if } \text{SoC} > \text{safe\_return\_threshold}(\text{distance\_to\_base}) \\
-100 & \text{if } \text{SoC} \leq \text{safe\_return\_threshold} \\
-1000 & \text{if collision or battery depleted}
\end{cases}
\end{align}
Safe return threshold is dynamically computed as the SoC required to reach base plus a 15\% safety margin, accounting for current resistance/assistance, wave drag, and battery degradation uncertainty.

Efficiency reward ($\lambda_{\text{eff}} = 0.3$) combines objectives:
\begin{align}
r^{\text{eff}}_t = \frac{\text{new\_coverage} \times \text{information\_gain}}{\text{energy\_consumed} + \epsilon}
\end{align}
rewarding high information density per Joule, where $\epsilon$ prevents division by zero during stationary information gathering.

Terminal rewards provide episode-level incentives. Mission completion ($\text{coverage} \geq 0.90$) gives $+500$ bonus. Safe return to base awards $+100$. Battery depletion yields $-1000$. Collision incurs $-500$. These terminal rewards shape overall episode trajectory toward successful missions.

\subsection{Network Architecture}

The policy network (approximately 250k parameters) accepts an 87D state vector and produces both discrete (4-way) and continuous (7D) action distributions. The architecture uses two hidden layers of 256 units with ReLU activation and dropout (rate 0.1) for regularization. The discrete head outputs logits for four motion modes via $\text{Linear}(256, 4) \to \text{softmax}$. The continuous head generates Gaussian parameters: mean via $\text{Linear}(256, 7) \to \text{Tanh}$ (bounded outputs), and log-std via $\text{Linear}(256, 7)$ clipped to $[\log(10^{-6}), \log(1.0)]$ for numerical stability. Actions are sampled as $a \sim \mathcal{N}(\mu, \sigma^2)$.

Twin critic networks (approximately 300k parameters each) accept 98-dimensional input concatenating the 87D state, 4D one-hot discrete action, and 7D continuous action. Each critic uses four linear layers (Linear(98, 256), Linear(256, 256), Linear(256, 128), Linear(128, 1)) with ReLU activation. Two independent networks $Q_1$ and $Q_2$ with different random initializations reduce overestimation bias that could lead to over-optimistic energy estimates.

Target networks are updated via exponential moving average with $\tau = 0.005$:
\begin{align}
\theta'_{\text{target}} \leftarrow (1 - \tau) \theta'_{\text{target}} + \tau \theta'
\end{align}

Temperature $\alpha$ (entropy coefficient) is learned via automatic entropy tuning to maintain target entropy $H_{\text{target}} = -7$ (negative of continuous action dimensionality):
\begin{align}
J(\alpha) = \mathbb{E}_{a \sim \pi}[-\alpha \log \pi(a|s) - \alpha H_{\text{target}}]
\end{align}

Network design prioritizes real-time deployment on embedded ASV platforms (NVIDIA Jetson AGX Orin) and resource-constrained edge devices. The lightweight architecture achieves approximately 12ms policy forward pass including state construction and action dispatch.

\subsection{Differentiation from Prior Work}

\begin{table*}[t]
\centering
\small
\begin{tabular}{lcccccc}
\hline
Method & Energy-Aware & Multi-Sensor & Water-Type & Joint Opt. & Safe Return & Real-Time \\
\hline
Frontier Exploration & No & No & No & No & No & Yes (fast) \\
Lawn-Mower CPP & No & No & No & No & No & Yes (fast) \\
Rule-Based Duty-Cycle & Partial & Yes & No & No & Yes & Yes \\
Energy-CPP (SAC) & Yes & No & No & Sequential & No & Yes ($\sim$15ms) \\
DDPG-ASV & Yes & No & No & No & No & Yes ($\sim$20ms) \\
ECMARL (Multi-Robot) & Yes & No & No & Yes & No & Medium ($\sim$50ms) \\
\textbf{Ours (Proposal 3)} & \textbf{Yes} & \textbf{Yes} & \textbf{Yes} & \textbf{Yes} & \textbf{Yes} & \textbf{Yes ($\sim$12ms)} \\
\hline
\end{tabular}
\vspace{0.3cm}
\caption{Comparison with existing autonomous surface vehicle and coverage planning methods. Energy-Aware = battery state tracking and consumption optimization. Multi-Sensor = joint duty-cycling of multiple sensors. Water-Type = explicit environment conditioning on water-body characteristics. Joint Opt. = simultaneous motion and sensing optimization. Safe Return = guaranteed energy reserve for return-to-base. Real-Time = inference latency within 12ms or better.}
\label{tab:comparison}
\end{table*}

Our approach introduces five key innovations beyond existing work.

First, joint motion-sensing optimization for the first time simultaneously optimizes ASV navigation thrust/steering and multi-sensor duty-cycling via a single unified RL policy. Prior work separates these concerns. Coverage path planning ignores sensors. Energy-aware path planning treats sensor scheduling as fixed or sequential. Energy-CPP from agricultural robotics optimizes coverage then energy in two stages, missing synergies such as reducing propulsion when valuable sensing is inactive.

Second, explicit water-body type conditioning enables adaptive strategies that exploit environment-specific opportunities. Rivers feature predictable currents (1--2 m/s steady flow) enabling energy-efficient drift-assisted exploration. Lakes provide stable conditions permitting aggressive coverage strategies. Coastal zones demand conservative wave-aware positioning. Open ocean requires maximum energy conservation. No prior ASV work conditions on water-body type, instead treating all environments uniformly.

Third, information-per-joule reward formulation directly optimizes the Pareto frontier between coverage and energy. Rather than weighted sum of coverage and energy as independent objectives, we compute $\frac{\text{area\_per\_sensor} \times \text{quality\_factor}}{\text{energy\_consumed}}$, explicitly trading off information gain against propulsion cost. This metric is novel to ASV exploration.

Fourth, guaranteed safe return constraint ensures sufficient energy reserve for return to base based on real-time distance and environmental resistance. Dynamic safe return threshold adapts to current conditions (current assistance/resistance), wave drag, and battery degradation. This safety guarantee reduces critical battery depletion events from 15\% (rule-based baselines) to less than 1\%.

Fifth, hierarchical action space with discrete mode selection plus continuous control provides interpretability and real-time efficiency while maintaining expressiveness. The four discrete modes (exploration, energy-saving, return-to-base, information-gathering) are transparent to human operators, enabling trust and debugging.
